\chapter{Applications Beyond Autonomous Vehicles}
\label{sec:applications}
% ============================================

% Springer Nature: each chapter should have a 10--15 line online abstract.
% This chapter abstract has not yet been written in the uploaded manuscript.
\abstract{\added{This chapter shows how the classical fusion toolbox of Chapter~\ref{sec:classicalSF} is assembled into complete monitoring systems outside autonomous driving. Two case studies are developed: maritime anomaly detection for the protection of critical infrastructure, and model-based monitoring of submarine cables with distributed acoustic sensing. In the first, AIS reports and coastal-radar plots are fused by a constant-velocity Kalman filter, and the innovation sequence is monitored with per-step, sliding-window and whiteness tests; a runnable example shows that a slow spoofing ramp and a coloured disturbance are exposed by different statistics. Behavioural features, route graphs and heterogeneous cues are then combined with evidence theory, whose conflict mass becomes an operator-relevant signal, and the resulting alerts are prioritised with AHP-weighted multi-criteria scores. In the second study, a damped tensioned-string model is reduced to a modal state space, DAS channels are described by a gauge-length measurement matrix, and a Kalman filter with windowed chi-square tests detects and localises an anchor strike along the cable. Worked numerical examples, figures generated by the companion repository, and end-of-chapter problems accompany both studies, and a closing paragraph draws the common pattern with the autonomous-driving pipeline: the same tools, applied to different sensors.}}
\suppressfloats[t]

% % Resolved: the leftover was replaced by the bridging paragraph below (blue).
The two case studies in this section are deliberately chosen to show how \emph{all} the classical ingredients of Chapter~\ref{sec:classicalSF} --- Kalman filtering and innovation monitoring, evidence-theoretic conflict handling, and MCDM-based prioritisation --- combine into complete fusion systems in a domain far from autonomous driving. They should be read as integrative worked examples of the survey's toolbox rather than as an exhaustive treatment of maritime sensing.

\section{Maritime Anomaly Detection and Infrastructure Protection}

Maritime anomaly detection is a representative application in which sensor fusion is needed because no single information source provides a complete and trustworthy maritime picture. The Automatic Identification System (AIS) provides vessel identity, position, speed, course, and other navigation information, but AIS messages can be missing, delayed, spoofed, or intentionally misleading. Coastal radar, electro-optical and infrared cameras, satellite imagery, synthetic-aperture radar, hydrophones, and distributed acoustic sensing can provide independent observations, but each modality has different coverage, resolution, latency, and false-alarm characteristics. Recent surveys of AIS anomaly detection highlight the diversity of abnormal behaviours that must be considered, including route deviation, loitering, unexpected stops, speed anomalies, encounter anomalies, and suspicious gaps in transmission \cite{wolsing_anomaly_2022}. This problem is becoming more relevant as maritime navigation systems become increasingly digital and interconnected, exposing the fused navigation pipeline itself to faults, spoofing, and cyber--physical disturbances \cite{oliva_innovation_2026}.
\begin{addedblock}
\subsection{Sources and Their Failure Modes}
AIS is a cooperative, self-reported source: each vessel broadcasts position reports (message types 1--3 for class~A transponders, type~18 for class~B) containing the Maritime Mobile Service Identity (MMSI), the GNSS-derived latitude and longitude, speed over ground (SOG), course over ground (COG), true heading, rate of turn and navigational status, together with static and voyage-related messages (type~5) that carry the vessel name, call sign, dimensions, type, draught and destination \cite{itu_m1371_2014}. The nominal reporting interval of a class~A transponder depends on the kinematic regime: 3\,min when anchored or moored, 10\,s between 0 and 14\,kn, $3\frac{1}{3}$\,s between 14 and 23\,kn, and 2\,s above 23\,kn, with the interval shortened when the vessel is changing course. Because the reports are unauthenticated and the receiving chain is passive, the AIS channel can fail in several characteristic ways: \emph{gaps} (transmitter switched off, VHF shadowing, or satellite-AIS revisit limits), \emph{identity inconsistencies} (an MMSI used by two hulls, or a static message that contradicts the observed size), \emph{kinematic inconsistencies} (reported SOG and COG incompatible with the displacement between consecutive positions), and \emph{spoofing}, in which forged or replayed messages place a vessel where it is not \cite{balduzzi_ais_2014}. The last two are the most insidious because a well-designed spoof respects the vessel dynamics and therefore passes the plausibility checks applied to a single source \cite{riveiro_review_2018}.

Non-cooperative sources are needed to close this gap. Coastal radar produces plots with range and bearing noise that, converted to Cartesian coordinates, grows with range; typical plot-to-plot standard deviations of a few tens of metres at 10--20\,km are common for X-band shore installations. Electro-optical and infrared cameras provide identity cues (silhouette, gear on deck, wake) over short ranges and in favourable visibility; satellite synthetic-aperture radar provides wide-area snapshots with revisit times of hours, which is why imagery is best used to confirm an existing track rather than to initiate one \cite{chen_sviadf_2025}. Distributed acoustic sensing on seabed fibres, discussed in Section~\ref{subsec:das-measurement-model}, adds a non-cooperative source that observes vessels only through the strain they induce in the cable, and whose spatial support is the cable itself. The design question addressed by the rest of this section is how to make these sources check one another.
\end{addedblock}


The practical difficulty is that suspicious behaviour is not always dramatic. A vessel may keep transmitting AIS in order to avoid immediate suspicion while still performing small route deviations, abnormal speed profiles, unusual loitering, or manoeuvres close to protected infrastructure. This is especially relevant in congested and geopolitically sensitive regions such as the Baltic Sea, where normal maritime traffic coexists with dense networks of submarine power and communication cables. A maritime fusion architecture should therefore combine trajectory analysis with contextual information: proximity to coastlines, shipping lanes, restricted areas, cable corridors, anchorages, recent incidents, weather, and vessel class.

At the trajectory level, fusion typically starts by constructing vessel tracks from AIS and radar measurements, associating visual or satellite detections with those tracks, and extracting behavioural features such as speed over ground, course change, distance from shipping lanes, proximity to restricted areas, encounter geometry, and consistency between reported identity and observed motion. Remote-sensing imagery and AIS can also be fused to identify small vessels or vessels whose reported AIS state is inconsistent with the observed scene \cite{chen_sviadf_2025}. Graph-based route models provide an additional layer of context: a recent Baltic Sea framework constructs a Probabilistic Roadmap (PRM) of feasible maritime corridors, projects AIS trajectories onto this graph, and derives interpretable features such as route inefficiency, deviation from learned corridors, and speed or dwell time near submarine cables \cite{oliva_tomei_prm_access_2026}. This illustrates a broader design principle: raw trajectories become more informative when fused with geospatial constraints and infrastructure maps.
\begin{addedblock}
\subsection{Track-Level Fusion of AIS and Radar}
\label{subsec:ais-radar-track-fusion}
The computations are carried out in a local east--north--up (ENU) frame with origin at the sensor site, so that the AIS latitude--longitude pairs and the radar range--bearing plots are first mapped to planar coordinates $(p_E,p_N)$. With the state $\mathbf{x}_k=[\,p_{E,k}\;\;p_{N,k}\;\;v_{E,k}\;\;v_{N,k}\,]^{\mathsf T}$, a merchant vessel in transit is well described by the constant-velocity (CV) model \eqref{eq:kf_linear_process} with
\begin{equation}
\mathbf{F}=
\begin{bmatrix}
\mathbf{I}_2 & T\mathbf{I}_2\\ \mathbf{0} & \mathbf{I}_2
\end{bmatrix},
\qquad
\mathbf{Q}=\sigma_a^{2}
\begin{bmatrix}
\tfrac{T^{4}}{4}\mathbf{I}_2 & \tfrac{T^{3}}{2}\mathbf{I}_2\\[2pt]
\tfrac{T^{3}}{2}\mathbf{I}_2 & T^{2}\mathbf{I}_2
\end{bmatrix},
\label{eq:ch07-cv-model}
\end{equation}
where $T$ is the report period and $\sigma_a$ bounds the unmodelled acceleration; manoeuvring vessels are better served by a coordinated-turn (CT) model in which the velocity vector rotates at a turn rate $\omega$, $v_{E,k+1}=v_{E,k}\cos\omega T-v_{N,k}\sin\omega T$ and $v_{N,k+1}=v_{E,k}\sin\omega T+v_{N,k}\cos\omega T$, either with $\omega$ appended to the state or through the multiple-model filters recalled in Chapter~\ref{sec:classicalSF} \cite{barshalom2001_estimation}. Both AIS and radar observe the position, $\mathbf{H}=[\,\mathbf{I}_2\;\;\mathbf{0}\,]$, with $\mathbf{R}_{\mathrm{AIS}}=\sigma_{\mathrm{AIS}}^{2}\mathbf{I}_2$ and $\mathbf{R}_{\mathrm{rad}}=\sigma_{\mathrm{rad}}^{2}\mathbf{I}_2$ (the latter obtained, in practice, by linearising the polar plot covariance at the current range). When the two sources arrive within the same period they are processed sequentially, the radar plot first and the AIS report second, each with its own innovation $\boldsymbol{\nu}_k$ and covariance $\mathbf{S}_k$ from \eqref{eq:kf_innovation}--\eqref{eq:kf_innovation_covariance}; the two innovations are mutually uncorrelated under a correct model, so their NIS values \eqref{eq:nis} add to a $\chi^{2}_{4}$ variable.

Association is performed before the update. A radar plot $j$ is assigned to a track only if it falls inside the gate \eqref{eq:mahalanobis_gate}, $\boldsymbol{\nu}_{k,j}^{\mathsf T}\mathbf{S}_{k,j}^{-1}\boldsymbol{\nu}_{k,j}\le\gamma$, with $\gamma=\chi^{2}_{2}(0.99)=9.21$ for a two-dimensional plot; AIS reports are associated by MMSI, but their position is still gated against the radar-supported prediction, and a report that fails the gate while its MMSI matches is precisely the signature of an identity or position inconsistency that the evidence layer of Section~\ref{subsec:maritime-evidence} will consume. In dense traffic, nearest-neighbour assignment is replaced by the probabilistic data association rules of Section~\ref{subsec:track-fusion}.

\subsection{Innovation-Based Monitoring of the Navigation Pipeline}
\label{subsec:maritime-innovation-monitoring}
Under a correct linear--Gaussian model the innovation sequence is zero-mean, Gaussian, white, and has covariance $\mathbf{S}_k$ \cite{barshalom2001_estimation,mehra_identification_1970}. Each of these properties yields a test, and the tests are not equivalent. The per-step NIS test flags sample $k$ when $d_k^{2}=\boldsymbol{\nu}_k^{\mathsf T}\mathbf{S}_k^{-1}\boldsymbol{\nu}_k>\chi^{2}_{n_y}(1-\alpha)$. The \emph{sliding-window} test accumulates $W$ consecutive values,
\begin{equation}
\epsilon_{k}^{W}=\sum_{i=k-W+1}^{k} d_i^{2}\;\sim\;\chi^{2}_{W n_y}
\quad\text{under } \mathcal{H}_0,
\qquad
\text{alarm if } \epsilon_{k}^{W}>\chi^{2}_{W n_y}(1-\alpha)
\;\text{ or }\;
\epsilon_{k}^{W}<\chi^{2}_{W n_y}(\alpha),
\label{eq:ch07-window-nis}
\end{equation}
where the two-sided form matters: an innovation that is systematically \emph{too small} is as much a symptom of model mismatch (an over-estimated $\mathbf{R}$, or a filter that is following a slowly varying disturbance) as one that is too large. The \emph{whiteness} test examines the sample autocorrelation of the whitened innovation components $\tilde{\nu}_{k,i}$, with $\tilde{\boldsymbol{\nu}}_k=\mathbf{S}_k^{-1/2}\boldsymbol{\nu}_k$,
\begin{equation}
\hat\rho_{i}(\ell)=\frac{\sum_{k=\ell+1}^{M}\tilde{\nu}_{k,i}\tilde{\nu}_{k-\ell,i}}{\sum_{k=1}^{M}\tilde{\nu}_{k,i}^{2}},
\qquad
Q_{\mathrm{LB}}=M(M+2)\sum_{i=1}^{n_y}\sum_{\ell=1}^{L}\frac{\hat\rho_{i}(\ell)^{2}}{M-\ell}\;\sim\;\chi^{2}_{n_y L},
\label{eq:ch07-ljung-box}
\end{equation}
computed over the last $M$ innovations and $L$ lags; this is the Ljung--Box portmanteau statistic \cite{ljung_box_1978}, applied component-wise and summed. A bias $\mathbf{b}$ in the innovation turns $d_k^{2}$ into a non-central $\chi^{2}$ variable with non-centrality $\lambda=\mathbf{b}^{\mathsf T}\mathbf{S}^{-1}\mathbf{b}$; the per-step test sees $\lambda$ once, whereas the windowed statistic sees $W\lambda$ against a threshold that grows only like $W n_y+O(\sqrt{W n_y})$, which is why a small persistent bias is detected by \eqref{eq:ch07-window-nis} long before any single sample exceeds its own threshold. A \emph{slow spoofing ramp} is the maritime instance of this situation. If a single sensor is spoofed and the ramp is dynamically plausible (a cross-track drift of a few centimetres per second squared is well within the $\sigma_a$ of a CV model), the filter that uses only that sensor simply follows the forged trajectory and its innovations remain white and consistent; the ramp becomes observable only in the innovation of an \emph{independent} sensor, and there it appears as a slowly growing bias, exactly the case for which the windowed test is designed. A \emph{coloured} disturbance (for instance a multipath or jamming-induced error that is correlated over several reports) leaves the marginal variance untouched and is invisible to the covariance tests, but it violates whiteness and is exposed by \eqref{eq:ch07-ljung-box} \cite{oliva_innovation_2026}. These three tests are complementary rather than redundant, as the following example shows.

\begin{example}[Innovation monitoring of a fused AIS--radar track.]
\label{ex:ch07-ais-innovation}
A vessel sails at $12$\,kn ($6.17$\,m/s) on course $060^{\circ}$ for 30\,min. A coastal radar provides a plot every $T=10$\,s with $\sigma_{\mathrm{rad}}=30$\,m per axis, and the class~A transponder reports its GNSS position at the same period with $\sigma_{\mathrm{AIS}}=10$\,m. The truth evolves according to \eqref{eq:ch07-cv-model} with $\sigma_a=0.01$\,m/s$^{2}$; the filter uses $\sigma_a=0.02$\,m/s$^{2}$ and processes the radar plot and the AIS report sequentially, so that the stacked per-step NIS is $\chi^{2}_{4}$ under $\mathcal{H}_0$. Three regimes are simulated: nominal operation (0--600\,s); a coloured AIS disturbance (600--1200\,s) in which the AIS noise becomes a first-order autoregressive process $e_k=\phi e_{k-1}+\sqrt{1-\phi^{2}}\,\sigma_{\mathrm{AIS}}w_k$ with $\phi=0.85$ and the \emph{same} marginal variance as before; and a slow spoofing ramp (from 1200\,s to the end) in which the reported AIS position drifts to starboard with a constant acceleration $a_s=0.001$\,m/s$^{2}$, reaching $180$\,m after 10\,min. All alarms use $\alpha=0.01$: the per-step threshold is $\chi^{2}_{4}(0.99)=13.28$; the window test uses $W=12$ (2\,min), hence $\chi^{2}_{48}(0.99)=73.68$; the Ljung--Box test uses $M=30$ innovations and $L=5$ lags, hence $\chi^{2}_{20}(0.99)=37.57$; and the AIS-only two-sided window test uses the $\chi^{2}_{24}$ interval $[9.89,\,45.56]$. In steady state $\sqrt{S_{11}}$ is $31.3$\,m for the radar and $13.2$\,m for the AIS update, and the position gains are $0.083$ and $0.429$: the fused track is driven mainly by AIS, which is what makes the AIS channel attractive to an adversary.

The outcome is summarised in Fig.~\ref{fig:ch07-ais-innovation}. In the nominal regime the mean NIS is $3.76$ (radar $1.64$, AIS $2.13$) and none of the tests fires. During the coloured disturbance the per-step and the stacked window tests never fire; the AIS innovations are in fact \emph{too small} (mean NIS $0.43$ instead of $2$) because the filter, which weights AIS heavily, follows the slowly varying error, and the root-mean-square position error grows from $7.5$ to $13.5$\,m. The whitened AIS innovations have a lag-one autocorrelation of $0.49$ (east) and $0.26$ (north); the Ljung--Box statistic exceeds its threshold in $9\,\%$ of the windows, first at $680$\,s, and the two-sided AIS-only window statistic falls below its lower bound in $94\,\%$ of the windows (median $5.1$ against an expectation of $24$), first at $690$\,s, i.e.\ $90$\,s after the onset. During the spoofing ramp the mean NIS rises to $8.91$ and, as anticipated, the excess sits in the \emph{radar} innovation ($6.22$ against $2$): the state follows the forged AIS and the trusted radar disagrees. The first alarms occur at $1580$\,s for the Ljung--Box test (offset $72$\,m), $1610$\,s for the window test ($84$\,m) and $1620$\,s for the per-step test ($88$\,m), after which the window test is active in $33\,\%$ of the samples, the whiteness test in $26\,\%$ and the per-step test in $21\,\%$; at the end of the scenario the fused track is $167$\,m from the true position. No test detects the ramp while the offset is below about $70$\,m, which is roughly two radar standard deviations: a spoof that is slow enough is indistinguishable from a manoeuvre until an independent source accumulates evidence against it. The script \repopath{companion\_\allowbreak repo/\allowbreak examples/\allowbreak ch07\_\allowbreak ais\_\allowbreak innovation\_\allowbreak monitoring/\allowbreak } reproduces every number quoted here.
\end{example}

\begin{figure}[!t]
\centering
\resizebox{0.99\columnwidth}{!}{
\begin{tikzpicture}
\begin{groupplot}[
    group style={group size=1 by 3, vertical sep=22mm},
    width=0.95\columnwidth,
    scale only axis,
    tick label style={font=\scriptsize},
    label style={font=\scriptsize},
    legend style={font=\scriptsize, draw=none, fill=white, fill opacity=0.85, text opacity=1},
    grid=major,
    grid style={line width=0.2pt, draw=gray!25},
    axis line style={line width=0.5pt},
    tick align=outside,
    title style={font=\scriptsize, yshift=0.5mm},
]
\nextgroupplot[
    height=3.0cm,
    xlabel={east $p_E$ [km]},
    ylabel={north $p_N$ [km]},
    xmin=-0.5, xmax=10.5, ymin=-0.5, ymax=6.5,
    legend pos=north west, legend cell align=left, legend columns=3,
    title={(a) fused AIS--radar track in the local ENU frame},
    x filter/.code={\pgfmathparse{#1/1000}\pgfmathresult},
    y filter/.code={\pgfmathparse{#1/1000}\pgfmathresult},
]
\addplot[only marks, mark=*, mark size=0.5pt, draw=MeasGrey, fill=MeasGrey]
    table[x=e_rad, y=n_rad, col sep=comma] {data/ch07/ais_innovation_monitoring.csv};
\addlegendentry{radar plots}
\addplot[only marks, mark=x, mark size=1.2pt, draw=ManoeuvreOrange]
    table[x=e_meas, y=n_meas, col sep=comma] {data/ch07/ais_innovation_monitoring.csv};
\addlegendentry{AIS reports}
\addplot[TruthBlack, line width=1.0pt]
    table[x=e_true, y=n_true, col sep=comma] {data/ch07/ais_innovation_monitoring.csv};
\addlegendentry{ground truth}
\addplot[EkfBlue, line width=0.9pt, densely dashed]
    table[x=e_est, y=n_est, col sep=comma] {data/ch07/ais_innovation_monitoring.csv};
\addlegendentry{KF}
\nextgroupplot[
    height=3.0cm,
    xlabel={time $t$ [s]},
    ylabel={$\epsilon_k^{W}$, $d_k^{2}$},
    xmin=0, xmax=1800, ymin=0, ymax=300,
    legend pos=north west, legend cell align=left, legend columns=2,
    title={(b) per-step NIS (grey) and windowed statistic ($W=12$) against $\chi^{2}_{48}(0.99)=73.7$},
]
\addplot[draw=none, fill=ManoeuvreOrange, fill opacity=0.35, forget plot]
    coordinates {(1200,0) (1800,0) (1800,300) (1200,300)} \closedcycle;
\addplot[draw=none, fill=BandGrey, forget plot]
    coordinates {(600,0) (1200,0) (1200,300) (600,300)} \closedcycle;
\addplot[draw=gray!70, line width=0.5pt, densely dotted, forget plot]
    coordinates {(0,73.68) (1800,73.68)};
\addplot[MeasGrey, line width=0.4pt]
    table[x=t, y=nis, col sep=comma] {data/ch07/ais_innovation_monitoring.csv};
\addlegendentry{$d_k^{2}$}
\addplot[EkfBlue, line width=0.7pt]
    table[x=t, y=nis_win, col sep=comma] {data/ch07/ais_innovation_monitoring.csv};
\addlegendentry{$\epsilon_k^{W}$}
\node[anchor=north, font=\scriptsize, text=black!70] at (axis cs:900,290) {coloured AIS noise};
\node[anchor=north, font=\scriptsize, text=black!70] at (axis cs:1500,290) {spoofing ramp};
\nextgroupplot[
    height=3.0cm,
    xlabel={time $t$ [s]},
    ylabel={$Q_{\mathrm{LB}}$, $\epsilon^{W}_{k,\mathrm{AIS}}$},
    xmin=0, xmax=1800, ymin=0, ymax=120,
    legend pos=north west, legend cell align=left, legend columns=2,
    title={(c) Ljung--Box whiteness statistic and AIS-only window statistic with its $\chi^{2}_{24}$ interval},
]
\addplot[draw=none, fill=ManoeuvreOrange, fill opacity=0.35, forget plot]
    coordinates {(1200,0) (1800,0) (1800,120) (1200,120)} \closedcycle;
\addplot[draw=none, fill=BandGrey, forget plot]
    coordinates {(0,9.89) (1800,9.89) (1800,45.56) (0,45.56)} \closedcycle;
\addplot[draw=gray!70, line width=0.5pt, densely dotted, forget plot]
    coordinates {(0,37.57) (1800,37.57)};
\addplot[EkfBlue, line width=0.7pt]
    table[x=t, y=lb, col sep=comma] {data/ch07/ais_innovation_monitoring.csv};
\addlegendentry{$Q_{\mathrm{LB}}$}
\addplot[TruthBlack, line width=0.6pt, densely dashed]
    table[x=t, y=nis_win_ais, col sep=comma] {data/ch07/ais_innovation_monitoring.csv};
\addlegendentry{$\epsilon^{W}_{k,\mathrm{AIS}}$}
\end{groupplot}
\end{tikzpicture}
}
\caption{Innovation monitoring of a fused AIS--radar track (Example~\ref{ex:ch07-ais-innovation}). (a)~Radar plots, AIS reports, ground truth and Kalman estimate; the AIS reports drift to starboard during the spoofing ramp. (b)~Per-step NIS and sliding-window statistic: the coloured AIS noise (grey band) is invisible to both, the spoofing ramp (orange) is detected once the forged position is about two radar standard deviations away from the radar-supported estimate. (c)~Ljung--Box statistic (threshold dotted) and AIS-only window statistic: the coloured disturbance drives the latter below the lower bound of its $\chi^{2}_{24}$ interval (grey band) because the filter follows the correlated error, and intermittently violates whiteness.}
\label{fig:ch07-ais-innovation}
\end{figure}

\subsection{Behavioural Features and Route Models}
\label{subsec:maritime-features}
The track-level estimate feeds a small set of interpretable features that are computed per report or per window. With the filtered velocity $(\hat v_{E,k},\hat v_{N,k})$, speed over ground and course over ground are $\mathrm{SOG}_k=\sqrt{\hat v_{E,k}^{2}+\hat v_{N,k}^{2}}$ and $\mathrm{COG}_k=\operatorname{atan2}(\hat v_{E,k},\hat v_{N,k})$, and the course-change rate is the wrapped difference $\dot\psi_k=\operatorname{wrap}(\mathrm{COG}_k-\mathrm{COG}_{k-1})/T$. The consistency between the reported and the estimated kinematics, $|\mathrm{SOG}^{\mathrm{AIS}}_k-\mathrm{SOG}_k|$ and $|\operatorname{wrap}(\mathrm{COG}^{\mathrm{AIS}}_k-\mathrm{COG}_k)|$, is itself a feature. Given a corridor graph with edges $\mathcal{E}$ obtained from historical traffic or from a probabilistic roadmap \cite{pallotta_tread_2013,oliva_tomei_prm_access_2026}, the distance to the nearest corridor is $d^{\mathrm{cor}}_k=\min_{e\in\mathcal{E}}\operatorname{dist}(\hat{\mathbf{p}}_k,e)$, and the route inefficiency over a voyage segment is the ratio between the travelled length and the shortest-path length on the graph. For a cable polyline $\mathcal{L}$, the distance $d^{\mathrm{cab}}_k=\operatorname{dist}(\hat{\mathbf{p}}_k,\mathcal{L})$ and the dwell time
\begin{equation}
\tau^{\mathrm{cab}}_k=T\sum_{i\le k}\mathbb{1}\!\left[d^{\mathrm{cab}}_i<d_0\;\wedge\;\mathrm{SOG}_i<s_0\right],
\label{eq:ch07-dwell}
\end{equation}
with, typically, $d_0=500$\,m and $s_0=2$\,kn, capture the behaviour that precedes anchoring or trawling over an asset. The surveys of \cite{riveiro_review_2018,wolsing_anomaly_2022} catalogue many further features; the ones above are those that map most directly onto the hypotheses of the evidence layer.
\end{addedblock}


Maritime anomaly detection should also include the estimation layer. In navigation and monitoring systems, GNSS, inertial sensors, radar, and AIS may be fused through Kalman-filter-based estimators. Some cyber--physical anomalies or stealthy disturbances may first appear as abnormal innovations before they become visible as large trajectory deviations. Innovation-based monitoring can therefore test whether the residual sequence remains zero-mean, covariance-consistent, Gaussian, and temporally white; such diagnostics can expose coloured spoofing or disturbance patterns that may bypass a simple \(\chi^2\) covariance check \cite{oliva_innovation_2026}. This connects ship anomaly detection directly to the sensor-fusion pipeline rather than treating it only as offline post-processing of AIS tracks.

Evidence theory is well suited to this setting because maritime anomalies are often ambiguous. A vessel moving slowly near a cable route may be fishing, drifting because of a mechanical problem, performing maintenance, or preparing an illegal anchoring activity. A frame of discernment can include hypotheses such as \emph{normal transit}, \emph{loitering}, \emph{route deviation}, \emph{AIS inconsistency}, \emph{collision risk}, \emph{infrastructure threat}, and \emph{unknown}. AIS-derived behaviour, radar tracks, camera detections, satellite observations, innovation-test alarms, and acoustic or cable-sensing cues can each provide BBAs over this frame. If several sources support the same abnormal hypothesis, the fused belief increases. If AIS reports normal transit while radar, imagery, or cable sensing indicates inconsistent behaviour, the conflict mass becomes an operator-relevant warning rather than being hidden inside a single confidence score.
\begin{addedblock}
\subsection{Evidence-Theoretic Fusion of Heterogeneous Cues}
\label{subsec:maritime-evidence}
\begin{example}[Four cues over a three-hypothesis frame.]
\label{ex:ch07-maritime-evidence}
Consider the reduced frame $\Omega=\{N,L,T\}$, where $N$ denotes normal transit, $L$ loitering or fishing, and $T$ a threat to the infrastructure (anchoring or tampering over a cable). Four sources provide BBAs about a vessel that reports ``underway using engine'' at 1.5\,kn, 300\,m from a cable corridor:
\begin{equation}
\small\setlength{\arraycolsep}{3pt}
\begin{array}{@{}l|ccccc@{}}
 & \{N\} & \{L\} & \{T\} & \{L,T\} & \Omega \\
\hline
m_{\mathrm{AIS}} & 0.60 & 0 & 0 & 0 & 0.40 \\
m_{\mathrm{rad}} & 0 & 0 & 0.20 & 0.50 & 0.30 \\
m_{\mathrm{img}} & 0 & 0.40 & 0 & 0.30 & 0.30 \\
m_{\mathrm{cab}} & 0 & 0 & 0.50 & 0.20 & 0.30
\end{array}
\label{eq:ch07-maritime-bbas}
\end{equation}
Here $m_{\mathrm{AIS}}$ reflects the declared transit and a plausible SOG, $m_{\mathrm{rad}}$ a radar plot 400\,m from the AIS position and a slow track, $m_{\mathrm{img}}$ an image showing gear deployed over the side, and $m_{\mathrm{cab}}$ a windowed-NIS alarm on the nearby cable; the masses on $\Omega$ encode the discounting of each source by its reliability, as in Chapter~\ref{sec:classicalSF}. Combining AIS and radar alone with Dempster's rule \eqref{eq:dempster} gives a conflict $K=0.60\times(0.20+0.50)=0.42$ and the fused belief $\mathrm{Bel}(T)=0.138$, $\mathrm{Pl}(T)=0.690$: the two sources disagree about the vessel's behaviour, and the conflict is the first warning. Adding imagery ($K=0.272$ with the running result) and then the cable cue ($K=0.241$) gives
\begin{equation}
\begin{aligned}
&m(\{T\})=0.450,\quad m(\{L\})=0.200,\quad m(\{L,T\})=0.266,\\
&m(\{N\})=0.051,\quad m(\Omega)=0.034,
\end{aligned}
\end{equation}
hence $\mathrm{Bel}(T)=0.450$, $\mathrm{Pl}(T)=0.750$, $\mathrm{Bel}(\{L,T\})=0.916$, and, by the pignistic transformation \eqref{eq:pignistic-transform}, $\mathrm{BetP}(N)=0.062$, $\mathrm{BetP}(L)=0.344$, $\mathrm{BetP}(T)=0.594$. The unnormalised Smets rule \eqref{eq:smets} applied to all four sources leaves a total mass $m(\emptyset)=0.680$ on the empty set, which is large: most of the product mass is in disagreement, essentially because the AIS declaration is incompatible with everything else. An operator therefore receives two numbers, a threat probability of about $0.6$ and a conflict of $0.68$, and the second is what justifies tasking a patrol or a camera rather than acting on the first alone. The companion script \repopath{companion\_\allowbreak repo/\allowbreak examples/\allowbreak ch07\_\allowbreak evidence\_\allowbreak mcdm/\allowbreak } prints these values.
\end{example}
\end{addedblock}


MCDM can complement this process by ranking alerts and response actions. Criteria may include distance to protected infrastructure, confidence of track association, deviation from learned maritime corridors, innovation-test alarms, severity of the suspected anomaly, persistence over time, weather and visibility, vessel type, traffic density, and expected operator workload. AHP-style pairwise comparisons can encode expert preferences from coast guards, port authorities, or infrastructure operators, while real-time sensor indicators provide the criterion scores. The resulting MCDM score can prioritise which anomalies require immediate attention, which sensors should be tasked for confirmation, or which vessels should be monitored more closely.
\begin{addedblock}
\subsection{Alert Prioritisation by MCDM}
\label{subsec:maritime-mcdm}
\begin{example}[Ranking four alerts under five AHP-weighted criteria.]
\label{ex:ch07-maritime-mcdm}
Suppose that, within one surveillance cycle, the evidence layer has produced four alerts: $a_1$, a vessel slowing down near a cable; $a_2$, a vessel leaving its corridor; $a_3$, an AIS--radar mismatch; and $a_4$, a vessel loitering in a designated anchorage. The criteria are $c_1$ proximity to protected infrastructure, $c_2$ confidence of the track association, $c_3$ deviation from the learned corridors, $c_4$ significance of the innovation-test alarm and $c_5$ persistence over time, all normalised so that $q_{ik}\in[0,1]$ increases with urgency:
\begin{equation}
\mathbf{Q}=
\begin{bmatrix}
0.90 & 0.80 & 0.40 & 0.60 & 0.70\\
0.30 & 0.90 & 0.90 & 0.20 & 0.50\\
0.60 & 0.50 & 0.30 & 0.95 & 0.40\\
0.20 & 0.95 & 0.10 & 0.10 & 0.90
\end{bmatrix}.
\end{equation}
An infrastructure operator provides the complete pairwise comparison matrix
\begin{equation}
\mathbf{A}=
\begin{bmatrix}
1 & 3 & 2 & 2 & 4\\
1/3 & 1 & 1/2 & 1/2 & 2\\
1/2 & 2 & 1 & 1 & 3\\
1/2 & 2 & 1 & 1 & 3\\
1/4 & 1/2 & 1/3 & 1/3 & 1
\end{bmatrix},
\end{equation}
whose principal eigenvector gives $\mathbf{w}=[\,0.376\;\;0.121\;\;0.215\;\;0.215\;\;0.074\,]^{\mathsf T}$ with $\lambda_{\max}=5.033$, consistency index $0.008$ and consistency ratio $0.007$, well below Saaty's $0.1$ limit \cite{saaty_analytic_1987}; the same weights are recovered by the logarithmic least-squares problem \eqref{eq:lls-ahp} on the complete graph. The additive scores \eqref{eq:mcdm-score} are $S(a_1)=0.701$, $S(a_3)=0.584$, $S(a_2)=0.495$, $S(a_4)=0.299$, so the slow vessel near the cable is served first. The ranking is, however, sensitive to how much the operator trusts the innovation layer: increasing $w_4$ while rescaling the other weights proportionally, $a_3$ overtakes $a_1$ at $w_4=0.412$, about twice its nominal value, whereas $a_4$ never reaches the top for any $w_4\in[0.05,0.6]$. Reporting such break-even weights alongside the ranking is a cheap robustness check of the kind advocated in Chapter~\ref{sec:classicalSF}.
\end{example}
\end{addedblock}


The main challenges are data quality, class imbalance, interpretability, and validation. True anomalous events are rare and often poorly labelled, while normal maritime traffic varies strongly with geography, season, weather, and port operations. Synthetic anomalies, incident-based case studies, and infrastructure-aware features can help evaluate candidate methods, but broader external validation remains necessary \cite{oliva_tomei_prm_access_2026}. For operational deployment, a transparent combination of track-level fusion, innovation monitoring, evidence-theoretic uncertainty handling, and MCDM-based prioritisation can be more useful than a black-box anomaly score alone.

\section{Sensor Fusion for Submarine Cable Monitoring}

Submarine power and telecommunication cables are critical infrastructures for digital connectivity, energy interconnection, offshore wind farms, and security-sensitive maritime operations \cite{carter_submarine_2009,kavanagh_achieving_2025}. They are exposed to accidental and intentional hazards such as anchor dragging, bottom trawling, seabed instability, fatigue, construction activity, and deliberate tampering. Anchor-related events are particularly important, with the risk depending on anchor geometry, burial depth, seabed properties, and vessel drift \cite{yoon_safety_2013,osthoff_anchor_soil_2017,li_anchor_risk_2025}. Monitoring is difficult because the relevant events are spatially distributed, the assets extend over tens or hundreds of kilometres, and damaging interactions may be masked by waves, currents, seabed coupling, and ordinary vessel activity.

Distributed acoustic sensing (DAS) turns an optical fibre into a dense array of strain or strain-rate measurements by interrogating backscattered light along the cable \cite{hartog_intro_2017,lu_distributed_2019,lindsey_fiberoptic_2021}. In marine environments, this is attractive because seafloor instrumentation is difficult and expensive to deploy, while existing or dedicated fibre-optic cables may cover long distances with metre-scale spatial sampling. Long-term subsea deployments show that DAS can support real-time detection of bottom-trawl fishing, ship anchoring, and seismic events over cable spans of operational length \cite{waagaard_longterm_2022}, and related demonstrations cover seismic, hydroacoustic, and power-cable strain sensing \cite{williams_distributed_2019,matsumoto_hydroacoustic_2021,masoudi_dynamic_2018}.
\begin{addedblock}
\subsection{Distributed Acoustic Sensing as a Measurement Model}
\label{subsec:das-measurement-model}
A phase-sensitive optical time-domain reflectometer ($\phi$-OTDR) launches pulses into the fibre and measures the optical phase of the Rayleigh backscatter as a function of the two-way travel time, hence of the position $x$ along the fibre. The phase difference between two points separated by the gauge length $L_g$ is proportional to the average axial strain over that interval, so that the $j$-th channel, centred at $x_j$, delivers
\begin{equation}
\varepsilon_j(t)=\frac{\kappa_j}{L_g}\left[w'(x_j+\tfrac{L_g}{2},t)-w'(x_j-\tfrac{L_g}{2},t)\right]+r_j(t),
\qquad j=1,\ldots,n_y,
\label{eq:ch07-das-channel}
\end{equation}
where $w'$ denotes the displacement of the cable projected on the fibre axis (for the transverse motion of a tensioned cable, the relevant quantity is the axial strain induced by bending and stretching of the armour, which we take proportional to $\partial w/\partial x$), $\kappa_j\in(0,1]$ is a channel-dependent coupling factor that accounts for the transfer of seabed motion to the fibre through the cable structure and for the burial or free-span condition of the segment, and $r_j$ is the interrogator noise, typically in the nano- to microstrain range depending on $L_g$, pulse energy and averaging \cite{hartog_intro_2017,lindsey_fiberoptic_2021}. Many interrogators output strain \emph{rate}, which is the time derivative of \eqref{eq:ch07-das-channel} and is handled by augmenting the state with the modal velocities, as done below. The channel spacing (metres to tens of metres) is independent of $L_g$ (typically 5--20\,m), so neighbouring channels are correlated whenever the spacing is smaller than the gauge length; the examples below use a spacing larger than $L_g$ so that $\mathbf{R}$ can be taken diagonal. The coupling factors are the DAS counterpart of the sensor-reliability indicators of Chapter~\ref{sec:classicalSF}: they enter $\mathbf{H}$ multiplicatively, and an incorrect $\kappa_j$ produces a channel whose innovations are systematically inconsistent, which the residual pattern of Section~\ref{subsec:cable-kf-monitoring} distinguishes from a mechanical event by its persistence and by its confinement to a single channel.
\end{addedblock}


For underwater-cable monitoring, DAS can support several fusion tasks. First, it can provide spatially distributed acoustic or deformation evidence of nearby vessels, anchors, trawling activity, sediment motion, or cable vibration. Second, it can be fused with AIS, radar, bathymetry, weather, ocean-current estimates, and maintenance records to interpret the cause of the observed signal. Third, it can support infrastructure protection by detecting events close to telecom cables, power cables, pipelines, offshore wind farms, or port approaches. Experimental work using an optical telecom cable offshore Toulon showed that ship-radiated noise could be detected and, in favourable conditions, used for trajectory estimation \cite{rivet_ship_2021}, while more recent Applied Ocean Research work studies DAS-based vessel monitoring using submarine fibre-optic cables in the presence of growing offshore infrastructure and vessel traffic \cite{paap_vessel_2025}. For power cables, \(\phi\)-OTDR experiments have also been used to distinguish calm sea, anchor drop, and anchor drag signatures \cite{lv_otdr_2018}.

The fusion problem is challenging because DAS measurements are not equivalent to conventional hydrophone measurements. The response depends on cable coupling to the seabed, cable orientation relative to the incoming wavefield, gauge length, interrogation settings, local bathymetry, and the propagation path. Some cable segments may be well coupled to the seafloor, while suspended or poorly coupled sections may respond differently. A robust fusion system should therefore include a cable-state model and should not treat all DAS channels as equally reliable.

Model-based monitoring can help bridge the gap between raw distributed measurements and actionable alarms. Cable dynamics can be approximated through physically grounded tensioned-structure models, with finite-dimensional state-space approximations providing a natural interface to filtering and diagnosis \cite{irvine_linear_1974,baudouin_cable_pde_2018}. 
\begin{addedblock}
\subsection{A Physics-Based Cable Model}
\label{subsec:cable-model}
A segment of length $L$ between two points where the cable is effectively pinned (burial transitions, supports, or the ends of a free span) is modelled as a damped tensioned string,
\begin{equation}
\rho\,\frac{\partial^{2} w}{\partial t^{2}}+c_d\,\frac{\partial w}{\partial t}-\mathcal{T}\,\frac{\partial^{2} w}{\partial x^{2}}=f(x,t),
\qquad w(0,t)=w(L,t)=0,
\label{eq:ch07-string-pde}
\end{equation}
where $w(x,t)$ is the transverse displacement, $\rho$ the mass per unit length, $c_d$ a viscous damping coefficient that lumps hydrodynamic and seabed dissipation, $\mathcal{T}$ the tension and $f$ the distributed load. Bending stiffness can be added as a term $EI\,\partial^{4}w/\partial x^{4}$ at the price of a different family of mode shapes \cite{irvine_linear_1974}; for the long spans of interest, tension dominates. Expanding $w(x,t)=\sum_{n=1}^{n_m} q_n(t)\varphi_n(x)$ in the mass-normalised mode shapes $\varphi_n(x)=\sqrt{2/L}\,\sin(n\pi x/L)$ and projecting \eqref{eq:ch07-string-pde} on each mode yields $n_m$ decoupled oscillators,
\begin{equation}
\ddot q_n+2\zeta_n\omega_n\dot q_n+\omega_n^{2}q_n=\frac{1}{\rho}\,f_n(t),
\qquad
\omega_n=\frac{n\pi}{L}\sqrt{\frac{\mathcal{T}}{\rho}},
\qquad
f_n(t)=\int_0^{L} f(x,t)\varphi_n(x)\,\mathrm{d}x,
\label{eq:ch07-modal}
\end{equation}
with $\zeta_n=c_d/(2\rho\omega_n)$. The state $\mathbf{x}=[\,\mathbf{q}^{\mathsf T}\;\;\dot{\mathbf{q}}^{\mathsf T}\,]^{\mathsf T}\in\mathbb{R}^{2n_m}$ obeys $\dot{\mathbf{x}}=\mathbf{A}\mathbf{x}+\mathbf{B}\mathbf{f}$ with
\begin{equation}
\mathbf{A}=
\begin{bmatrix}
\mathbf{0} & \mathbf{I}\\ -\boldsymbol{\Omega}^{2} & -2\mathbf{Z}\boldsymbol{\Omega}
\end{bmatrix},
\qquad
\mathbf{B}=\frac{1}{\rho}
\begin{bmatrix}
\mathbf{0}\\ \mathbf{I}
\end{bmatrix},
\qquad
\boldsymbol{\Omega}=\operatorname{diag}(\omega_n),\;\;\mathbf{Z}=\operatorname{diag}(\zeta_n),
\label{eq:ch07-modal-ss}
\end{equation}
and the environmental excitation (currents, surface waves transmitted through the water column, micro-seismic noise) is represented by a white modal force with spectral density $\mathbf{Q}_c$. Discretising with period $\Delta t$ gives $\mathbf{F}=\mathrm{e}^{\mathbf{A}\Delta t}$ and the process covariance $\mathbf{Q}=\int_0^{\Delta t}\mathrm{e}^{\mathbf{A}s}\mathbf{B}\mathbf{Q}_c\mathbf{B}^{\mathsf T}\mathrm{e}^{\mathbf{A}^{\mathsf T}s}\,\mathrm{d}s$, both computed exactly from a single matrix exponential of the block matrix $\bigl[\begin{smallmatrix}-\mathbf{A} & \mathbf{B}\mathbf{Q}_c\mathbf{B}^{\mathsf T}\\ \mathbf{0} & \mathbf{A}^{\mathsf T}\end{smallmatrix}\bigr]\Delta t$. Substituting the modal expansion in \eqref{eq:ch07-das-channel} shows that the DAS measurement is linear in the state,
\begin{equation}
\mathbf{y}_k=\mathbf{H}\mathbf{x}_k+\mathbf{r}_k,
\qquad
H_{jn}=\frac{\kappa_j}{L_g}\left[\varphi_n\!\left(x_j+\tfrac{L_g}{2}\right)-\varphi_n\!\left(x_j-\tfrac{L_g}{2}\right)\right],
\quad
H_{j,n_m+n}=0,
\label{eq:ch07-das-H}
\end{equation}
i.e.\ $\mathbf{H}$ samples the finite differences of the mode shapes at the channel positions, and a channel centred on a point where $\varphi_n'$ vanishes is blind to mode $n$. For strain-rate interrogators the non-zero block moves to the velocity columns. The pair \eqref{eq:ch07-modal-ss}--\eqref{eq:ch07-das-H} is exactly the linear model \eqref{eq:kf_linear_process}--\eqref{eq:kf_linear_measurement} of Chapter~\ref{sec:classicalSF}, so the Kalman recursion \eqref{eq:kf_state_prediction}--\eqref{eq:kf_cov_update} applies without modification.
\end{addedblock}


In a recent paper~\cite{fioravanti_cable_med_2026} this direction is followed by modelling a submarine cable as a damped tensioned structure, obtaining a finite-dimensional state-space approximation, and applying Kalman filtering to distributed deformation measurements. Under nominal conditions, the filter innovations should remain statistically consistent with the model and the assumed environmental excitation. Localised mechanical interactions, such as an anchor strike or tampering event, appear as systematic deviations in the innovation sequence and can be detected using sliding-window statistical tests. This approach complements purely data-driven DAS processing because it introduces physical interpretability and separates stochastic environmental forcing from persistent structural perturbations.
\begin{addedblock}
\subsection{Kalman Filtering and Innovation Monitoring on Distributed Measurements}
\label{subsec:cable-kf-monitoring}
With $n_y$ channels the per-step NIS is $\chi^{2}_{n_y}$ and the windowed statistic \eqref{eq:ch07-window-nis} is $\chi^{2}_{W n_y}$ under $\mathcal{H}_0$. A mechanical interaction enters \eqref{eq:ch07-modal} as an unmodelled modal force $f_n(t)=F(t)\varphi_n(x_e)$ concentrated at the contact point $x_e$; its effect on the innovations is a bias that is spread over the modes and, through $\mathbf{H}$, over the channels. Detection and localisation are therefore separated. Detection uses the scalar statistic $\epsilon_k^{W}$, which aggregates all channels and is insensitive to where the event is. Localisation uses the channel-wise normalised residuals $\tilde\nu_{k,j}=\nu_{k,j}/\sqrt{S_{k,jj}}$ accumulated over the alarm window $\mathcal{W}$,
\begin{equation}
E_j=\max\Bigl\{0,\;\sum_{k\in\mathcal{W}}\tilde\nu_{k,j}^{2}-|\mathcal{W}|\Bigr\},
\qquad
\hat x_e=\frac{\sum_j x_j E_j}{\sum_j E_j},
\label{eq:ch07-residual-profile}
\end{equation}
where the subtraction removes the $\mathcal{H}_0$ expectation of one per sample and channel, and the centroid $\hat x_e$ is a model-free estimate of the contact position. Because the channels measure strain, the residual pattern of a point load is antisymmetric about $x_e$: the profile $E_j$ has a minimum \emph{at} the contact point flanked by two maxima, and the resolution is bounded by the modal truncation, roughly $L/(2n_m)$. A model-based alternative is to augment the state with the unknown force at a grid of candidate positions and to compare the likelihoods, in the spirit of the multiple-model filters of Chapter~\ref{sec:classicalSF}; the profile \eqref{eq:ch07-residual-profile} is cheaper and is sufficient to decide which segment to inspect.

\begin{example}[Anchor strike on a 1-km span.]
\label{ex:ch07-cable-das}
Let $L=1000$\,m, $\sqrt{\mathcal{T}/\rho}=50$\,m/s, $\rho=20$\,kg/m and $\zeta_n=0.05$, with $n_m=5$ modes at $0.025$, $0.050$, $0.075$, $0.100$ and $0.125$\,Hz; the model is sampled at $\Delta t=0.5$\,s for 140\,s. Sixteen DAS channels are centred every 60\,m between $x=50$ and $950$\,m with $L_g=10$\,m; two channels ($x=290$ and $350$\,m) lie on a free span and have $\kappa_j=0.5$, the others $\kappa_j=1$, and $\sigma_r=0.05$\,$\mu\varepsilon$. The white environmental forcing produces a nominal strain of $1.04$\,$\mu\varepsilon$ RMS. At $t_e=60$\,s an anchor strikes the cable at $x_e=640$\,m: an impulse of $0.7$\,N\,s is applied to the modal velocities, $\Delta\dot q_n=(I_0/\rho)\varphi_n(x_e)$, followed by a $0.1$\,N drag force sustained for 15\,s; the peak strain reaches $6.4$\,$\mu\varepsilon$. The filter knows $\mathbf{F}$, $\mathbf{Q}$, $\mathbf{H}$ (including the coupling factors) and $\mathbf{R}$ but not the strike. With $\alpha=0.01$ the per-step threshold is $\chi^{2}_{16}(0.99)=32.0$ and the window test with $W=8$ (4\,s) uses $\chi^{2}_{128}(0.99)=168.1$.

Figure~\ref{fig:ch07-cable-das} shows the result. Before the strike the mean NIS is $15.3$ against an expectation of $16$ and the window test never fires. At the strike sample the per-step NIS is $20.5$, \emph{below} its threshold: the impulse first changes the modal velocities, which the strain channels do not observe directly, so no single sample is grossly inconsistent. The windowed statistic, however, crosses its threshold at $t=60.5$\,s, one sample after the event, peaks at $245$, and stays above threshold for $4$\,s while the filter re-converges. Over this alarm window the residual-energy profile \eqref{eq:ch07-residual-profile} has its two largest values at $x=770$ and $830$\,m and a secondary maximum at $530$\,m, with a dip at $650$\,m, the channel nearest the contact; the centroid is $\hat x_e=696$\,m against a true $640$\,m, an error of $56$\,m for a modal resolution of $L/(2n_m)=100$\,m. The free-span channels, whose halved coupling is known to the filter, contribute nothing anomalous. Problem~\ref{prob:ch07:7} explores how the localisation error changes with the number of modes and with an unknown coupling factor; the script is in \repopath{companion\_\allowbreak repo/\allowbreak examples/\allowbreak ch07\_\allowbreak cable\_\allowbreak das\_\allowbreak kf/\allowbreak }.
\end{example}

\begin{figure}[!t]
\centering
\resizebox{0.99\columnwidth}{!}{
\begin{tikzpicture}
\begin{groupplot}[
    group style={group size=1 by 2, vertical sep=20mm},
    width=0.95\columnwidth,
    scale only axis,
    tick label style={font=\scriptsize},
    label style={font=\scriptsize},
    legend style={font=\scriptsize, draw=none, fill=white, fill opacity=0.85, text opacity=1},
    grid=major,
    grid style={line width=0.2pt, draw=gray!25},
    axis line style={line width=0.5pt},
    tick align=outside,
    title style={font=\scriptsize, yshift=0.5mm},
]
\nextgroupplot[
    height=3.0cm,
    xlabel={time $t$ [s]},
    ylabel={$\epsilon_k^{W}$, $d_k^{2}$},
    xmin=0, xmax=140, ymin=0, ymax=260,
    legend pos=north west, legend cell align=left, legend columns=2,
    title={(a) per-step NIS and windowed statistic ($W=8$) against $\chi^{2}_{128}(0.99)=168.1$},
]
\addplot[draw=none, fill=ManoeuvreOrange, fill opacity=0.45, forget plot]
    coordinates {(60,0) (75,0) (75,260) (60,260)} \closedcycle;
\addplot[draw=gray!70, line width=0.5pt, densely dotted, forget plot]
    coordinates {(0,168.1) (140,168.1)};
\addplot[draw=gray!70, line width=0.5pt, densely dotted, forget plot]
    coordinates {(0,128) (140,128)};
\addplot[MeasGrey, line width=0.4pt]
    table[x=t, y=nis, col sep=comma] {data/ch07/cable_das_timeseries.csv};
\addlegendentry{$d_k^{2}$}
\addplot[EkfBlue, line width=0.7pt]
    table[x=t, y=nis_win, col sep=comma] {data/ch07/cable_das_timeseries.csv};
\addlegendentry{$\epsilon_k^{W}$}
\node[anchor=west, font=\scriptsize, text=black!70] at (axis cs:78,230) {anchor strike and drag};
\nextgroupplot[
    height=3.0cm,
    xlabel={position along the cable $x$ [m]},
    ylabel={$E_j$},
    xmin=0, xmax=1000, ymin=0, ymax=80,
    legend pos=north west, legend cell align=left,
    title={(b) residual-energy profile over the alarm window; true contact point dashed},
]
\addplot[draw=TruthBlack, dashed, line width=0.6pt, forget plot]
    coordinates {(640,0) (640,80)};
\addplot[draw=EkfBlue, densely dotted, line width=0.6pt, forget plot]
    coordinates {(696,0) (696,80)};
\addplot[ybar, bar width=14pt, fill=EkfBlue!60, draw=EkfBlue]
    table[x=x, y=E, col sep=comma] {data/ch07/cable_das_profile.csv};
\addlegendentry{$E_j$ per channel}
\node[anchor=west, font=\scriptsize, text=black!70] at (axis cs:60,70) {free-span channels ($\kappa_j=0.5$) at 290 and 350\,m};
\node[anchor=east, font=\scriptsize, text=black!70] at (axis cs:630,70) {$x_e=640$\,m};
\node[anchor=west, font=\scriptsize, text=EkfBlue] at (axis cs:705,55) {$\hat x_e=696$\,m};
\end{groupplot}
\end{tikzpicture}
}
\caption{Model-based DAS monitoring of a cable span (Example~\ref{ex:ch07-cable-das}). (a)~The windowed chi-square statistic detects the anchor strike one sample after the event although no single per-step NIS exceeds $\chi^{2}_{16}(0.99)=32$; the lower dotted line is the expectation $Wn_y=128$. (b)~Normalised residual energy per channel over the alarm window: the strain response to a point load is antisymmetric about the contact point, so the profile dips at the nearest channel and peaks on either side; its centroid (dotted) lies within the modal resolution of the true position (dashed).}
\label{fig:ch07-cable-das}
\end{figure}
\end{addedblock}


Evidence theory can be used to represent uncertainty in event interpretation. For a cable segment and time window, the hypotheses may include \emph{passing vessel}, \emph{anchoring or trawling risk}, \emph{seismic event}, \emph{oceanographic background}, \emph{cable-coupling artefact}, \emph{mechanical tampering}, and \emph{unknown}. DAS spectral features, distributed deformation residuals, AIS proximity, radar tracks, hydrophone data, bathymetry, and environmental models can each assign belief masses. High conflict, such as a strong cable disturbance without a nearby AIS vessel, may indicate an unreported vessel, missing data, a non-vessel geophysical source, or a localised mechanical interaction requiring inspection.
\begin{addedblock}
\subsection{Evidence and Decision Layers}
\label{subsec:cable-evidence-decision}
\begin{example}[Interpreting the alarm of Example~\ref{ex:ch07-cable-das}.]
\label{ex:ch07-cable-evidence}
Take $\Omega_c=\{V,A,S,C\}$ for passing vessel, anchor or trawl contact, seismic event and coupling artefact. The DAS spectrum of the alarm window is broadband at low frequency, $m_{\mathrm{spec}}(\{V,A\})=0.50$, $m_{\mathrm{spec}}(\{S\})=0.10$, $m_{\mathrm{spec}}(\Omega_c)=0.40$; the residual profile is localised and persists for 4\,s, $m_{\mathrm{res}}(\{A,C\})=0.60$, $m_{\mathrm{res}}(\Omega_c)=0.40$; the AIS picture shows a vessel within 300\,m at 1.5\,kn, $m_{\mathrm{AIS}}(\{V,A\})=0.70$, $m_{\mathrm{AIS}}(\Omega_c)=0.30$; and the regional seismic network reports no event, $m_{\mathrm{seis}}(\{V,A,C\})=0.80$, $m_{\mathrm{seis}}(\Omega_c)=0.20$. Dempster's rule \eqref{eq:dempster} produces pairwise conflicts of $0.060$, $0.030$ and $0.011$ as the sources are added, and the fused BBA $m(\{A\})=0.519$, $m(\{A,V\})=0.346$, $m(\{A,C\})=0.080$, $m(\{A,C,V\})=0.043$, $m(\Omega_c)=0.011$, $m(\{S\})=0.003$, whose pignistic probabilities are $\mathrm{BetP}(A)=0.748$, $\mathrm{BetP}(V)=0.190$, $\mathrm{BetP}(C)=0.057$, $\mathrm{BetP}(S)=0.005$. Contrary to the maritime case of Example~\ref{ex:ch07-maritime-evidence}, the sources agree and the conflict is small: the appropriate reaction is to raise an anchor-contact alert and to pass it to the prioritisation layer, where it competes with other alerts under the criteria listed below exactly as in Example~\ref{ex:ch07-maritime-mcdm}, with cable criticality and inspection cost replacing corridor deviation and persistence.
\end{example}
\end{addedblock}


MCDM is useful for prioritising alerts and sensor resources. Criteria may include signal-to-noise ratio, distance to the cable, estimated vessel size, consistency with AIS, innovation-test significance, persistence of the event, cable criticality, environmental conditions, and cost of dispatching inspection assets. The same framework can rank candidate explanations or decide whether to request additional confirmation from patrol vessels, satellite imagery, or nearby fixed sensors. In this way, DAS-based cable monitoring becomes a full sensor-fusion problem: dense acoustic or deformation sensing provides the raw evidence, physical models and filters structure the residuals, and contextual maritime data and decision logic turn that evidence into actionable situation awareness.

\begin{addedblock}
Taken together, the two case studies exhibit the same pattern as the autonomous-driving pipeline of Chapter~\ref{sec:application}. A physically motivated linear model (constant velocity for a vessel, modal dynamics for a cable) is paired with a measurement model that encodes the sensor geometry (an ENU projection, a gauge-length difference of mode shapes); a Kalman filter fuses the sources and, as a by-product, delivers an innovation sequence whose consistency, size and whiteness are monitored with the per-step, windowed and portmanteau tests of \eqref{eq:ch07-window-nis}--\eqref{eq:ch07-ljung-box}; the alarms, together with cues that have no state-space description (imagery, route graphs, spectral signatures), are expressed as BBAs over a small frame and combined with \eqref{eq:dempster} or \eqref{eq:smets}, the conflict being reported rather than normalised away; and the surviving alerts are ranked with \eqref{eq:mcdm-score} under weights elicited through \eqref{eq:lls-ahp}. What changes from the road to the sea is the sensor suite, the time scale (seconds for a lane change, minutes for a spoofing ramp, hours for a loitering pattern) and the cost of a wrong decision, not the tools.

\problemsection
\begin{prob}\label{prob:ch07:1}
Two sources report on the frame $\Omega=\{N,L,T\}$ of Example~\ref{ex:ch07-maritime-evidence}: an AIS cue $m_1(\{N\})=0.7$, $m_1(\Omega)=0.3$ and a radar cue $m_2(\{T\})=0.4$, $m_2(\{L,T\})=0.4$, $m_2(\Omega)=0.2$. Compute by hand the conflict $K$, the Dempster-combined BBA \eqref{eq:dempster}, $\mathrm{Bel}(T)$, $\mathrm{Pl}(T)$ and the pignistic probabilities \eqref{eq:pignistic-transform}. Then discount the AIS source with reliability $\alpha_1=0.5$ and repeat; comment on how the conflict and $\mathrm{BetP}(T)$ change.
\end{prob}
\begin{prob}\label{prob:ch07:2}
A fused AIS--radar filter processes $n_y=4$ scalar innovation components per 10-s period. Derive the distribution of the sliding-window statistic \eqref{eq:ch07-window-nis} for $W=12$ under $\mathcal{H}_0$ and compute its one-sided $99\,\%$ threshold and its two-sided $99\,\%$ interval (use a $\chi^{2}$ table or the Wilson--Hilferty approximation $\chi^{2}_{\nu}(p)\approx\nu\bigl[1-\tfrac{2}{9\nu}+z_p\sqrt{\tfrac{2}{9\nu}}\bigr]^{3}$). If the window is advanced by one sample at a time, explain why the number of alarms per hour under $\mathcal{H}_0$ cannot be read directly from $\alpha$, and give an upper bound on the expected number of alarm \emph{samples} per hour.
\end{prob}
\begin{prob}\label{prob:ch07:3}
Let a constant bias $\mathbf{b}$ affect the innovation of a sensor with innovation covariance $\mathbf{S}$. Show that the per-step NIS is non-central $\chi^{2}_{n_y}$ with non-centrality $\lambda=\mathbf{b}^{\mathsf T}\mathbf{S}^{-1}\mathbf{b}$ and that the windowed statistic is non-central $\chi^{2}_{Wn_y}$ with non-centrality $W\lambda$. For the radar of Example~\ref{ex:ch07-ais-innovation} ($\sqrt{S_{11}}=31.3$\,m, $n_y=2$, $W=12$), compute the mean of each statistic for a $50$-m bias and compare it with the respective $99\,\%$ thresholds; explain why the windowed test is the more sensitive one for a slow ramp.
\end{prob}
\begin{prob}\label{prob:ch07:4}
Using the score matrix $\mathbf{Q}$ and the weights $\mathbf{w}$ of Example~\ref{ex:ch07-maritime-mcdm}, verify the four scores $S(a_i)$. Then find analytically the value of the proximity weight $w_1$ (rescaling the other weights proportionally) at which $a_3$ would overtake $a_1$, or show that no such value exists in $[0,1)$.
\end{prob}
\begin{prob}\label{prob:ch07:5}
For the string of Section~\ref{subsec:cable-model} with $L=1000$\,m and $L_g=10$\,m, write the entries $H_{jn}$ of \eqref{eq:ch07-das-H} for a channel centred at $x_j=250$\,m and $n=1,2,3,4$ (take $\kappa_j=1$). Show that the channel is blind to mode $n=2$ and explain, using the small-gauge approximation $H_{jn}\approx\kappa_j\varphi_n'(x_j)$, which modes a channel at $x_j=500$\,m cannot observe. What does this imply for the choice of channel positions when the filter is used for localisation?
\end{prob}
\begin{prob}\label{prob:ch07:6}
(Computational.) Starting from the script of Example~\ref{ex:ch07-ais-innovation}, vary the spoofing acceleration $a_s\in\{0.0005,0.001,0.002,0.004\}$\,m/s$^{2}$ and the window length $W\in\{6,12,24\}$ and tabulate, for each pair, the detection delay of the per-step and windowed tests and the AIS offset at detection. Then repeat the coloured-disturbance regime with $\phi\in\{0.5,0.7,0.85,0.95\}$ and report the fraction of Ljung--Box alarms. See \repopath{companion\_\allowbreak repo/\allowbreak exercises/\allowbreak ch07/\allowbreak }.
\end{prob}
\begin{prob}\label{prob:ch07:7}
(Computational.) Starting from the script of Example~\ref{ex:ch07-cable-das}, move the strike to $x_e\in\{200,400,640,800\}$\,m and use $n_m\in\{3,5,8\}$ modes in both the simulator and the filter; report the detection delay and the localisation error of the centroid \eqref{eq:ch07-residual-profile}. Then keep $n_m=5$ and let the filter assume $\kappa_j=1$ for the two free-span channels whose true coupling is $0.5$: describe how the residual profile distinguishes this coupling artefact from a mechanical event. See \repopath{companion\_\allowbreak repo/\allowbreak exercises/\allowbreak ch07/\allowbreak }.
\end{prob}
\end{addedblock}
